\chapterimage{figs/chapterhead.png}
\chapter{The Main Graphical Application}
\label{chap:main-graphical-application}

This chapter documents the current \texttt{genesys-gui} surface as it exists
in the \texttt{WorkInProgress} branch at commit
\texttt{7cc6458da93c5099d7f1f5dda7535f5a761999ea}. It is grounded in the
present Qt window layout and in the controllers that now split the main window
into narrower responsibilities.

\noindent\textbf{Objective.} Describe the real top-level GUI, the visible
panels, and the user-facing actions that anchor day-to-day model work.

\noindent\textbf{Audience.} Users who start with the graphical interface and
need to understand where to load models, edit them, inspect traces, and launch
simulation runs.

\noindent\textbf{Scope.} Cover the main window layout, menus, toolbars, model
trees, property editor, console, result panes, simulation controls, tool
actions, and the safe shutdown path.

\noindent\textbf{Status.} Partially validated by code inspection and by the
current branch structure; visual screenshots are intentionally represented by
compilable placeholders until the final images are produced.

The GUI is the main composition root for the desktop workflow. It wires the
model lifecycle, scene interaction, trace rendering, property editing, plugin
catalog updates, and simulation event handling through dedicated controllers.
That makes the window easier to document in slices: the frame itself, the
navigation widgets, the editable scene, the trace and results panes, and the
simulation configuration path.

The rest of this chapter follows that order through
Section~\ref{sec:main-graphical-application-layout},
Section~\ref{sec:main-graphical-application-menus},
Section~\ref{sec:main-graphical-application-trees},
Section~\ref{sec:main-graphical-application-scene},
Section~\ref{sec:main-graphical-application-trace},
Section~\ref{sec:main-graphical-application-simulation},
Section~\ref{sec:main-graphical-application-actions}, and
Section~\ref{sec:main-graphical-application-closure}.

\section{Main window layout}
\label{sec:main-graphical-application-layout}

The central window is built by \texttt{mainwindow.ui} and initialized by
\texttt{mainwindow.cpp}. The root widget owns the standard Qt menus, the
toolbars, the central tab sets, the dock widgets, and the scene/view pair that
renders the current model. The central surface is split across:
\begin{itemize}
\item \texttt{tabWidgetCentral};
\item \texttt{tabWidgetModel};
\item \texttt{tabWidgetSimulation};
\item \texttt{tabWidgetReports};
\item \texttt{dockWidgetPropertyEditor};
\item \texttt{dockWidgetConsole};
\item the graphical view backed by \texttt{ModelGraphicsScene}.
\end{itemize}

The layout is intentionally split into a model-editing center and supporting
inspection panes. The editor area keeps the scene, text model, simulation
trace tabs, and report tabs visible as separate work surfaces instead of
collapsing everything into one monolithic canvas. That separation matches the
current controller split and makes it easier to explain which actions are
always available and which are disabled while a simulation is running.

% FIGURE-SPEC-BEGIN
% ID: fig:main-graphical-window-annotated
% Source of truth:
%   - source/applications/gui/genesys/mainwindow.ui
%   - source/applications/gui/genesys/mainwindow.h
%   - source/applications/gui/genesys/mainwindow.cpp
% Purpose:
%   Show the main Qt window with its dock widgets, tab widgets, scene view, and primary action areas.
% Required visual content:
%   - full main window frame
%   - central model/simulation/report tabs
%   - property editor and console docks
%   - annotated scene/view region
% Accessibility description:
%   A labeled overview of the main GUI frame so the reader can map the rest of the chapter to the window.
% Validation:
%   The screenshot must match the current WorkInProgress branch and the visible widgets named in mainwindow.ui.
% FIGURE-SPEC-END
\begin{figure}[htbp]
\centering
\figureplaceholder{figs/generated/main-graphical-window-annotated}
\caption{Annotated main GUI window showing the primary workspace and supporting panes.}
\label{fig:main-graphical-window-annotated}
\end{figure}

Figure~\ref{fig:main-graphical-window-annotated} is the anchor for the rest of
this chapter. It should be read together with the menus in
Figure~\ref{fig:main-graphical-menus} and the scene/editor panes in
Figures~\ref{fig:main-graphical-component-tree},
\ref{fig:main-graphical-property-editor}, and
\ref{fig:main-graphical-graphics-area}.

\section{Menus and toolbars}
\label{sec:main-graphical-application-menus}

The top-level menus are the normal entry points for model lifecycle, editing,
simulation, tools, and about/help actions. The current menu surface includes
Model, Edit, View, Simulation, Tools, and About. The Model menu owns new,
open, recent, save, close, information, check, and model-navigation actions.
The Edit menu exposes undo, redo, find, replace, clipboard, grouping, and
ungrouping actions. The View and drawing toolbars control zoom, grid/ruler
visibility, guides, snapping, alignment, and graphical selection behavior.

The toolbar set is arranged as a compact action strip around the main window.
The current code adds the model, edit, view, arrange, draw, simulation,
graphical-model, animate, and about toolbars to the top and left areas, then
locks them down to the top area at runtime so the workspace stays predictable.
That is the correct reading of the GUI: the toolbars are a command surface, not
independent windows.

Several actions are intentionally state-gated. When no model is open, the GUI
disables most editing and simulation controls. When a simulation is active, it
locks the edit, view, drawing, and property-editing paths so the runtime state
cannot be damaged by concurrent edits.

% FIGURE-SPEC-BEGIN
% ID: fig:main-graphical-menus
% Source of truth:
%   - source/applications/gui/genesys/mainwindow.ui
%   - source/applications/gui/genesys/mainwindow.cpp
% Purpose:
%   Show the menu bar and toolbar cluster with the current action groupings.
% Required visual content:
%   - Model/Edit/View/Simulation/Tools/About menus
%   - the top and left toolbar groups
%   - visible enabled/disabled states for representative actions
% Accessibility description:
%   A guide to where the user finds the main commands and how they are grouped.
% Validation:
%   The screenshot must reflect the current menu names and runtime gating behavior.
% FIGURE-SPEC-END
\begin{figure}[htbp]
\centering
\figureplaceholder{figs/generated/main-graphical-menus}
\caption{Menu bar and toolbar grouping in the main GUI.}
\label{fig:main-graphical-menus}
\end{figure}

\section{Model trees and property editor}
\label{sec:main-graphical-application-trees}

The left-hand navigation and the property editor are where the current model
becomes inspectable. The component tree lists model components and reflects the
scene selection. The data-definition tree lists the current model data
definitions. The plugin tree shows loaded plugins and their catalog entries.
The property editor on the right mirrors the selected model object and is kept
in sync with scene selection changes and tree selection changes.

The important behavior here is synchronization. Selecting an item in the tree
or in the scene should drive the property editor to the same model object. The
code keeps that linkage in the composition root and in the property-editor
controller rather than leaving it implicit. That is why the chapter should
describe the trees as coordinated views of the model, not as independent
widgets.

% FIGURE-SPEC-BEGIN
% ID: fig:main-graphical-component-tree
% Source of truth:
%   - source/applications/gui/genesys/mainwindow.ui
%   - source/applications/gui/genesys/mainwindow.h
%   - source/applications/gui/genesys/controllers/ModelInspectorController.cpp
% Purpose:
%   Show the component tree and its relation to the model and scene selection.
% Required visual content:
%   - treeWidgetComponents
%   - treeWidgetDataDefnitions or adjacent data-definition tree
%   - treeWidget_Plugins if visible in the same workspace capture
% Accessibility description:
%   A structural view of the model browser area that explains how items are organized.
% Validation:
%   The screenshot must show actual model objects from the documented commit, not placeholders.
% FIGURE-SPEC-END
\begin{figure}[htbp]
\centering
\figureplaceholder{figs/generated/main-graphical-component-tree}
\caption{Model browsing area with the component and data-definition trees.}
\label{fig:main-graphical-component-tree}
\end{figure}

% FIGURE-SPEC-BEGIN
% ID: fig:main-graphical-property-editor
% Source of truth:
%   - source/applications/gui/genesys/mainwindow.ui
%   - source/applications/gui/genesys/mainwindow.cpp
%   - source/applications/gui/genesys/controllers/PropertyEditorController.cpp
% Purpose:
%   Show the docked property editor used to edit the selected model object.
% Required visual content:
%   - property editor dock
%   - representative editable fields
%   - visible selection context where practical
% Accessibility description:
%   A close-up of the inspector used to edit object properties from the scene or trees.
% Validation:
%   The screenshot must reflect the current selection-driven property browser behavior.
% FIGURE-SPEC-END
\begin{figure}[htbp]
\centering
\figureplaceholder{figs/generated/main-graphical-property-editor}
\caption{Docked property editor synchronized with the current selection.}
\label{fig:main-graphical-property-editor}
\end{figure}

\section{Scene and graphical editing}
\label{sec:main-graphical-application-scene}

The graphics view is the primary editing surface for model diagrams.
The backing scene is \texttt{ModelGraphicsScene}.
That scene receives the draw, arrange, zoom, grouping, connection, and
visibility commands exposed by the menus and toolbars. The visible editing
modes include the normal selection flow and the drawing and animation actions
for lines, rectangles, ellipses, polygons, text, counters, variables, queues,
stations, entities, events, attributes, statistics, plots, and simulated
time.

The chapter should also make clear that the scene is not just a canvas. It is
the coordination point for undo, graphical connections, selection, visibility
filters, and graphical simulation overlays. That is why the current GUI
maintains separate actions for grid, ruler, guides, snap, recursive visibility,
attached elements, and editable elements. Those options do not merely change
appearance; they change how the model is navigated and edited.

% FIGURE-SPEC-BEGIN
% ID: fig:main-graphical-graphics-area
% Source of truth:
%   - source/applications/gui/genesys/mainwindow.ui
%   - source/applications/gui/genesys/mainwindow.cpp
%   - source/applications/gui/genesys/controllers/SceneToolController.cpp
% Purpose:
%   Show the graphics view and the visible scene editing aids.
% Required visual content:
%   - model graphics area
%   - grid/ruler/guides/snap state where possible
%   - a representative model diagram or open scene
% Accessibility description:
%   The main interactive canvas for placing and connecting model components.
% Validation:
%   The screenshot must match the scene tools and visibility controls documented in this chapter.
% FIGURE-SPEC-END
\begin{figure}[htbp]
\centering
\figureplaceholder{figs/generated/main-graphical-graphics-area}
\caption{Graphical editing area and scene-level interaction aids.}
\label{fig:main-graphical-graphics-area}
\end{figure}

\section{Trace, console, and reports}
\label{sec:main-graphical-application-trace}

The lower and side output panes are where the GUI exposes runtime feedback.
The main console receives normal traces and errors. The simulation pane shows
simulation-specific traces. The reports pane shows report-oriented output and
the summary results that are emitted when the model or runtime suppresses a
full automatic report dialog.

This is implemented through the trace bridge controller instead of through
ad hoc widget writes in the window class. The practical effect for the reader
is that the console, simulation output, and reports output are separate sinks
with different meanings. A general trace does not prove a simulation completed,
and a report pane does not mean the model was scientifically validated. The
chapter should say that explicitly so the interface is not overread.

Figure~\ref{fig:main-graphical-trace-panel} captures the trace-oriented panes,
while Figure~\ref{fig:main-graphical-results-panel} captures the report and
result-oriented panes that appear after a run.

% FIGURE-SPEC-BEGIN
% ID: fig:main-graphical-trace-panel
% Source of truth:
%   - source/applications/gui/genesys/mainwindow.ui
%   - source/applications/gui/genesys/controllers/TraceConsoleController.cpp
%   - source/applications/gui/genesys/mainwindow.cpp
% Purpose:
%   Show the console and simulation trace panes with their live output roles.
% Required visual content:
%   - main console pane
%   - simulation trace pane
%   - report text pane when visible
% Accessibility description:
%   The runtime feedback region used for traces, errors, and simulation messages.
% Validation:
%   The screenshot must correspond to the live trace routing described in code.
% FIGURE-SPEC-END
\begin{figure}[htbp]
\centering
\figureplaceholder{figs/generated/main-graphical-trace-panel}
\caption{Console and trace panes used for runtime feedback.}
\label{fig:main-graphical-trace-panel}
\end{figure}

% FIGURE-SPEC-BEGIN
% ID: fig:main-graphical-results-panel
% Source of truth:
%   - source/applications/gui/genesys/mainwindow.cpp
%   - source/applications/gui/genesys/controllers/SimulationEventController.cpp
%   - source/tools/Statistics/SimulationResultsDataset.cpp
% Purpose:
%   Show the reports and results region after a simulation or replication completes.
% Required visual content:
%   - reports tab or results table
%   - representative summary output or plot region
%   - state that indicates finished or suppressed report display
% Accessibility description:
%   The area where a user inspects summaries, traces, and result-oriented views.
% Validation:
%   The screenshot must come from the current reporting behavior and not from mock data.
% FIGURE-SPEC-END
\begin{figure}[htbp]
\centering
\figureplaceholder{figs/generated/main-graphical-results-panel}
\caption{Results and reports area after execution.}
\label{fig:main-graphical-results-panel}
\end{figure}

\section{Simulation controls and configuration}
\label{sec:main-graphical-application-simulation}

The simulation menu and the dedicated toolbar expose start, stop, pause,
resume, step, and configuration actions. The configuration action opens the
simulation configuration dialog, which mirrors the live \texttt{ModelSimulation}
state and lets the user edit replication counts, time units, warm-up period,
termination condition, step-by-step execution, pause behavior, and local
parallelization settings. The dialog also carries a distributed-parallelization
placeholder button, which is explicitly documented in the code as a future
extension rather than as a finished backend integration.

This is the right place to mention the state gating again. The GUI enables the
simulation actions only when the current model and runtime state permit them.
That means the chapter should describe the controls as available commands with
preconditions, not as always-active buttons.

Figure~\ref{fig:main-graphical-simulation-config} shows the dialog surface the
user reaches from this command.

% FIGURE-SPEC-BEGIN
% ID: fig:main-graphical-simulation-config
% Source of truth:
%   - source/applications/gui/genesys/dialogs/dialogsimulationconfigure.ui
%   - source/applications/gui/genesys/dialogs/dialogsimulationconfigure.h
%   - source/applications/gui/genesys/dialogs/dialogsimulationconfigure.cpp
%   - source/applications/gui/genesys/controllers/ModelLifecycleController.cpp
% Purpose:
%   Show the simulation configuration dialog used by the main GUI.
% Required visual content:
%   - replication and time-unit fields
%   - warm-up and termination settings
%   - pause/step-by-step controls
%   - local parallelization controls
%   - distributed-parallelization placeholder area or button
% Accessibility description:
%   A configuration dialog for the runtime parameters that drive simulation execution.
% Validation:
%   The dialog capture must match the live fields loaded from ModelSimulation and ExperimentManager.
% FIGURE-SPEC-END
\begin{figure}[htbp]
\centering
\figureplaceholder{figs/generated/main-graphical-simulation-config}
\caption{Simulation configuration dialog opened from the main GUI.}
\label{fig:main-graphical-simulation-config}
\end{figure}

\section{Available, incomplete, and experimental actions}
\label{sec:main-graphical-application-actions}

The current GUI exposes a small number of actions that should be documented as
present but not necessarily mature. The visible example is the legacy
\texttt{ToolsExperimentation} slot, which remains in the class but does not
currently map to a live QAction in \texttt{mainwindow.ui}. That is a
compatibility artifact, not a user-facing feature. The chapter should call
such items out plainly so readers can distinguish implemented commands from
historical wiring and future scaffolding.

More generally, the GUI uses explicit state checks to keep incomplete or
dangerous actions disabled when they should not run. The documentation should
reflect that contract instead of implying that every listed slot is a polished
feature. This is especially important for the simulation toolbar, the plugin
manager entry point, and the graphical simulation toggle.

\section{Safe closure}
\label{sec:main-graphical-application-closure}

Closing the main window is not a trivial hide-and-exit operation. The current
window class coordinates pending-change checks, scene cleanup, trace shutdown,
signal disconnection, and model persistence prompts before allowing the
application to exit. In practice, that means the user can close the window
without losing the current model accidentally, but only after the current
dirty state has been reviewed.

The chapter should end by telling the user that the GUI is a coordinated
workspace, not a thin shell around a single canvas. The model tree, property
editor, scene, trace panes, and simulation dialog all participate in that
workspace and are kept consistent by the controller layer.
